




\begin{quote}\textbf{Résumé.} 
Ce volume est une \'edition recompos\'ee et annot\'ee du livre \og Rev\^etements \'Etales et Groupe Fondamental\fg, Lecture Notes in Mathematics, 224, Springer-Verlag, Berlin-Heidelberg-New York, 1971, par Alexander Grothendieck et al.\\
Le texte pr\'esente les fondements d'une th\'eorie du groupe fondamental en G\'eom\'etrie Alg\'ebrique, dans le point de vue \og kroneckerien\fg permettant de traiter sur le m\^eme pied le cas d'une vari\'et\'e alg\'ebrique au sens habituel, et celui d'un anneau des entiers d'un corps de nombres, par exemple.
\end{quote}













\cleardoublepage
\renewcommand{\baselinestretch}{1.05}\normalfont
\chapter*{Pr\'eface}
\thispagestyle{empty}
\vfill
Le pr\'esent texte est une \'edition recompos\'ee
et annot\'ee
du livre \og Rev\^etements
\'Etales et Groupe Fondamental\fg, Lecture Notes in Mathematics, 224,
Springer-Verlag, Berlin-Heidelberg-New York, 1971, par Alexander
Grothendieck et al.

La composition en \LaTeX 2e a \'et\'e r\'ealis\'ee par des volontaires, participant \`a un projet dirig\'e par Bas Edixhoven, dont on trouvera plus de d\'etails \`a l'adresse \url{http://www.math.leidenuniv.nl/~edix/}. La mise en page a \'et\'e termin\'ee par la \hbox{Société} math\'ematique de France.

Ceci est une \emph{version légèrement corrigée} du texte
original. Elle est \'edit\'ee par la \hbox{Société} math\'ematique de France (\url{http://smf.emath.fr/}) dans la s\'erie \og Documents \hbox{Mathématiques}\fg. Quelques mises \`a jour ont \'et\'e effectu\'ees par Michel Raynaud. Celles-ci sont d\'elimit\'ees par des crochets \lcrochetbf\,\rcrochetbf\ et indiqu\'ees par le symbole (MR) (remarques pages \pageref{X.2.14}, \pageref{XI.1.4}, \pageref{XII.5.6}, \pageref{XIII.2.13} et note \ref{III.6.6.p24} page \pageref{III.6.6.p24}). Afin d'uniformiser les notations, le corps r\'esiduel d'un point  $x$  est not\'e  $\kres(x)$,
et celui d'un anneau local  $A$  est not\'e  $\kres(A)$.\par
Il existe \'egalement une version \'electronique cens\'ee reproduire le texte original.

Les deux versions ont un seul fichier source commun, se trouvant sur
le serveur des archives arXiv.org e-Print \`a
\url{http://arxiv.org/}. Les diff\'erences entre les deux versions
sont document\'ees dans le fichier source.

L'ancienne num\'erotation des pages est incorpor\'ee dans la marge
(le nombre~$n$ indiquant le d\'ebut de la page~$n$).

\vfill
\newpage
The present text is a new
updated
edition of the book ``Rev\^etements \'Etales et
Groupe Fondamental'', Lecture Notes in Mathematics, 224,
Springer-Verlag, Berlin-Heidelberg-New York, 1971, by Alexander
Grothendieck et al.

Typesetting in \LaTeX 2e was done by volunteers, participating in a project directed by Bas Edixhoven, on which one can find more details at
\url{http://www.math.leidenuniv.nl/~edix/}. It has been completed by the Soci\'et\'e \hbox{mathématique} de France.

This is a slightly \emph{corrected version} of the original
text. It is published by the \hbox{Société} \hbox{mathématique} de France (\url{http://smf.emath.fr/}) as a volume of the series ``\hbox{Documents} Math\'ematiques''. Some updating remarks have been added by Michel Raynaud. These are bounded by brackets \lcrochetbf\,\rcrochetbf\ and indicated by the symbol (MR) (Remarks on pages \pageref{X.2.14}, \pageref{XI.1.4}, \pageref{XII.5.6}, \pageref{XIII.2.13} and footnote \ref{III.6.6.p24} on page \pageref{III.6.6.p24}). In order to unify notation, the residue field of a point $x$ is denoted by $\kres(x)$ and the residue field of a local ring $A$ is denoted by $\kres(A)$.\par
There also exists an electronic version that is meant to reproduce the original text.

Both versions are produced from the same source file, that can be
downloaded from the arXiv.org e-Print server at
\url{http://arxiv.org/}. All differences between the two versions are
documented in the source file.

The old page numbering is incorporated in the margin (the number~$n$
marking the beginning of page~$n$).

\vfill

\cleardoublepage
\renewcommand{\baselinestretch}{1.1}\normalfont
\chapter*{Introduction}
\thispagestyle{empty}
\label{I.0}


Dans la premi\`ere partie de cette introduction, nous donnons des
pr\'ecisions sur le contenu du pr\'esent volume; dans la
deuxi\`eme, sur l'ensemble du \og \emph{Séminaire de
Géométrie Algébrique du Bois-Marie}\fg, dont le pr\'esent
volume constitue le tome premier.

\let\oldthesubsection\thesubsection
\def\thesubsection{\arabic{subsection}}
\subsection{}
\label{I.introduction.1}
Le pr\'esent volume pr\'esente les fondements d'une
th\'eorie du groupe fondamental en G\'eom\'etrie Alg\'ebrique,
dans le point de vue \og kroneckerien\fg permettant de traiter sur le
m\^eme pied le cas d'une vari\'et\'e alg\'ebrique au sens
habituel, et celui d'un anneau des entiers d'un corps de nombres, par
exemple. Ce point de vue ne s'exprime d'une fa\c con satisfaisante
que dans le langage des sch\'emas, et nous utiliserons librement ce
langage, ainsi que les r\'esultats principaux expos\'es dans les
trois premiers chapitres des \emph{Éléments de Géométrie
Algébrique} de J.\, \textsc{Dieudonné} et A.\, \textsc{Grothendieck}, (cit\'e EGA dans
la suite). L'\'etude du pr\'esent volume du \og \emph{Séminaire
de Géométrie Algébrique du Bois-Marie}\fg ne demande pas
d'autres connaissances de la G\'eom\'etrie Alg\'ebrique, et peut
donc servir d'introduction aux techniques actuelles de
G\'eom\'etrie Alg\'ebrique, \`a un lecteur d\'esireux de se
familiariser avec ces techniques.

Les expos\'es \ref{I} \`a \ref{XI} de ce livre sont une reproduction
textuelle, pratiquement inchang\'ee, des notes
mim\'eographi\'ees du S\'eminaire oral, qui \'etaient
distribu\'ees par les soins de l'\emph{Institut des Hautes Études
Scientifiques\footnote{Ainsi que les notes des séminaires faisant
suite à celui-ci. Ce mode de distribution s'étant
avéré impraticable et insuffisant à la longue, tous les
\og \emph{Séminaire de Géométrie Algébrique du
Bois-Marie}\fg paraîtront désormais sous forme de livre comme
le présent volume.}}. Nous nous sommes born\'es \`a rajouter
quelques notes de bas de page au texte primitif, de corriger quelques
erreurs de frappe, et de faire un ajustage terminologique, le mot
\og morphisme simple\fg ayant notamment \'et\'e remplac\'e
entre-temps par celui de \og morphisme lisse\fg, qui ne pr\^ete pas aux
m\^emes confusions.

Les expos\'es \ref{I} \`a \ref{IV} pr\'esentent les notions locales de
morphisme \emph{étale} et de morphisme \emph{lisse}; ils
n'utilisent gu\`ere le langage des sch\'emas, expos\'e dans le
Chapitre~I des \emph{Éléments}\footnote{Une \'etude plus
compl\`ete est maintenant disponible dans les \emph{Éléments},
Chap~IV, \oldS\S 17 et 18.}. L'expos\'e~\ref{V} pr\'esente la description
axiomatique du groupe fondamental d'un sch\'ema, utile m\^eme dans
le cas classique o\`u ce sch\'ema se r\'eduit au spectre d'un
corps, o\`u on trouve une reformulation fort commode de la
th\'eorie de Galois habituelle. Les expos\'es \ref{VI} et \ref{VIII}
pr\'esentent la \emph{théorie de la descente}, qui a pris une
importance croissance en G\'eom\'etrie Alg\'ebrique dans ces
derni\`eres ann\'ees, et qui pourrait rendre des services
analogues en G\'eom\'etrie Analytique et en Topologie. Il convient
de signaler que l'expos\'e VII n'avait pas \'et\'e
r\'edig\'e, et sa substance se trouve
incorpor\'ee
dans un
travail de J.\, Giraud (M\'ethode de la Descente, Bull.\ Soc.\ Math.\
France, M\'emoire~2, 1964, viii + 150~p.). Dans l'expos\'e \ref{IX}, on
\'etudie plus sp\'ecifiquement la descente des morphismes
\'etales, obtenant une approche syst\'ematique pour des
th\'eor\`emes du type de \textsc{Van Kampen} pour le groupe fondamental,
qui apparaissent ici comme de simples traductions de th\'eor\`emes
de descente. Il s'agit essentiellement d'un proc\'ed\'e de calcul
du groupe fondamental d'un sch\'ema connexe~$X$, muni d'un morphisme
surjectif et propre, disons $X'\to X$, en termes des groupes
fondamentaux des composantes connexes de~$X'$ et des produits
fibr\'es~$X'\times_X X'$, $X'\times_X X'\times_X X'$, et des
homomorphismes induits entre ces groupes par les morphismes
simpliciaux canoniques entre les sch\'emas
pr\'ec\'edents. L'expos\'e~\ref{X} donne la th\'eorie de la
\emph{spécialisation du groupe fondamental}, pour un morphisme
propre et lisse, dont le r\'esultat le plus frappant consiste en la
d\'etermination (\`a peu de chose pr\`es) du groupe fondamental
d'une courbe alg\'ebrique lisse en caract\'eristique $p>0$,
gr\^ace au r\'esultat connu par voie transcendante en
caract\'eristique nulle. L'expos\'e~\ref{XI} donne quelques exemples et
compl\'ements, en explicitant notamment sous forme cohomologique la
th\'eorie des rev\^etements de \emph{\textsc{Kummer}}, et celle
d'\emph{\textsc{Artin}-\textsc{Schreier}}. Pour d'autres commentaires sur le texte,
voir l'\emph{Avertissement} \`a la version multigraphi\'ee, qui
fait suite \`a la pr\'esente Introduction.

Depuis la r\'edaction en 1961 du pr\'esent S\'eminaire a
\'et\'e d\'evelopp\'e, en collaboration par M.\, \textsc{Artin} et
moi-m\^eme, le langage de la \emph{topologie étale} et une
th\'eorie cohomologique correspondante, expos\'ee dans la partie
SGA~4 \og Cohomologie \'etale des sch\'emas\fg du \emph{Séminaire
de Géométrie Algébrique}, \`a para\^itre dans la
m\^eme s\'erie que le pr\'esent volume. Ce langage, et les
r\'esultats dont il dispose d\`es \`a pr\'esent, fournissent
un outil particuli\`erement souple pour l'\'etude du groupe
fondamental, permettant de mieux comprendre et de d\'epasser
certains des r\'esultats expos\'es ici. Il y aurait donc lieu de
reprendre enti\`erement la th\'eorie du groupe fondamental de ce
point de vue (tous les r\'esultats-clefs figurant en fait d\`es
\`a pr\'esent dans \loccit). C'est ce qui \'etait projet\'e
pour le chapitre des \emph{Éléments} consacr\'e au groupe
fondamental, qui devait contenir \'egalement plusieurs autres
d\'eveloppements qui n'ont pu trouver leur place ici (s'appuyant sur
la technique de r\'esolution des singularit\'es): calcul du
\og groupe fondamental local\fg d'un anneau local complet en termes d'une
r\'esolution des singularit\'es convenable de cet anneau, formules
de K\"unneth locales et globales pour le groupe fondamental sans
hypoth\`ese de propret\'e (\cf Exp.\, \ref{XIII}), les r\'esultats de
M.\, \textsc{Artin} sur la comparaison des groupes fondamentaux locaux d'un
anneau local hens\'elien excellent et de son compl\'et\'e
(SGA~4~XIX). Signalons \'egalement la n\'ecessit\'e de
d\'evelopper une th\'eorie du groupe fondamental d'un topos, qui
englobera \`a la fois la th\'eorie topologique ordinaire, sa
version semi-simpliciale, la variante \og profinie\fg
 d\'evelopp\'ee
dans l'expos\'e~\ref{V} du pr\'esent volume, et la variante
pro-discr\`ete un peu plus g\'en\'erale de SGA~3~X~7
(adapt\'ee au cas de sch\'emas non normaux et non unibranches). En
attendant une refonte d'ensemble de la th\'eorie dans cette optique,
l'expos\'e \ref{XIII} de Mme \textsc{Raynaud}, utilisant le langage et les
r\'esultats de SGA~4, est destin\'e \`a montrer le parti qu'on
peut tirer dans quelques questions typiques, en g\'en\'eralisant
notamment certains r\'esultats de l'expos\'e~\ref{X} \`a des
sch\'emas relatifs non propres. On y donne en particulier la
structure du groupe fondamental \og premier \`a $p$\fg d'une courbe
alg\'ebrique non compl\`ete en
caract\'eristique
quelconque (que j'avais
annonc\'ee
en 1959, mais dont une d\'emonstration n'avait pas
\'et\'e publi\'ee \`a ce jour).

Malgr\'e ces nombreuses lacunes et imperfections (d'autres diront:
\`a cause de ces lacunes et imperfections), je pense que le
pr\'esent volume pourra \^etre utile au lecteur qui d\'esire se
familiariser avec la th\'eorie du groupe fondamental, ainsi que
comme ouvrage de r\'ef\'erence, en attendant la r\'edaction et
la parution d'un texte \'echappant aux critiques que je viens
d'\'enum\'erer.

\subsection{}
\label{I.introduction.2}
Le pr\'esent volume constitue le tome 1 du
\og \emph{Séminaire de Géométrie Algébrique du
Bois-Marie}\fg, dont les volumes suivants sont pr\'evus pour
para\^itre dans la m\^eme s\'erie que celui-ci. Le but que se
propose le \emph{Séminaire}, parall\`element au trait\'e
\og \emph{Éléments de Géométrie Algébrique}\fg de J.\, \textsc{Dieudonné} et
A.\, \textsc{Grothendieck}, est de jeter les fondements de la G\'eom\'etrie
Alg\'ebrique, suivant les points de vue dans ce dernier ouvrage. La
r\'ef\'erence standard pour tous les volumes du
\emph{Séminaire} est constitu\'ee par les Chapitres~I, II, III
des \og \emph{Éléments de Géométrie Algébrique}\fg
(cit\'es EGA~I, II, III), et on suppose le lecteur en possession du
bagage d'alg\`ebre commutative et l'alg\`ebre homologique que ces
chapitres impliquent\footnote{Voir l'Introduction \`a EGA~I pour des
pr\'ecisions \`a ce sujet.}. De plus, dans chaque volume du
\emph{Séminaire} il sera r\'ef\'er\'e librement, dans le
mesure des besoins, \`a des volumes ant\'erieurs du m\^eme
\emph{Séminaire}, ou \`a d'autres chapitres publi\'es ou sur
le point de para\^itre des \og \'El\'ements\fg.

Chaque \emph{partie} du \emph{Séminaire} est centr\'ee sur un
sujet principal, indiqu\'e dans le titre du ou des volumes
correspondants; le s\'eminaire oral porte g\'en\'eralement sur
une ann\'ee acad\'emique, parfois plus. Les expos\'es \`a
l'int\'erieur de chaque partie du \emph{Séminaire} sont
g\'en\'eralement dans une d\'ependance logique \'etroite les
uns par rapport aux autres; par contre, les diff\'erentes parties du
\emph{Séminaire} sont dans une large mesure logiquement
ind\'ependants les uns par rapport aux autres. Ainsi, la partie
\og Sch\'emas en Groupes\fg est \`a peu pr\`es enti\`erement
ind\'ependante des deux parties du \emph{Séminaire} qui la
pr\'ec\`edent chronologiquement; par contre, elle fait un
fr\'equent appel aux r\'esultats de EGA~IV. Voici la liste des
parties du \emph{Séminaire} qui doivent para\^itre
prochainement (cit\'es SGA~1 \`a SGA~7 dans la suite):
\begin{enumerateb}
\item[SGA 1.] Rev\^etements \'etales et groupe fondamental, 1960
et 1961.
\item[SGA 2.] Cohomologie locale des faisceaux coh\'erents et
th\'eor\`emes de Lefschetz locaux et globaux, 1961/62.
\item[SGA 3.] Sch\'emas en groupes, 1963 et 1964 (3~volumes),
en~coll. avec M.\, \textsc{Demazure}.
\item[SGA 4.] Th\'eorie des topos et cohomologie \'etale des
sch\'emas, 1963/64 (3~volumes), (en~coll. avec M.\, \textsc{Artin} et
J.L.\, \textsc{Verdier}).
\item[SGA 5.] Cohomologie $l$-adique et fonctions~$L$, 1964 et 1965
(2~volumes).
\item[SGA 6.] Th\'eorie des intersections et th\'eor\`eme de
Riemann-Roch, 1966/67 (2~volumes) (en~coll. avec P.\, \textsc{Berthelot} et
L.\, \textsc{Illusie}).
\item[SGA 7.] Groupes de monodromie locale en g\'eom\'etrie
alg\'ebrique.
\end{enumerateb}

Trois parmi ces \emph{Séminaires} partiels ont \'et\'e
dirig\'es en \emph{collaboration} avec d'autres
math\'e\-maticiens, qui figureront comme co-auteurs sur la
couverture des volumes correspondants. Quant aux autres participants
actifs du \emph{Séminaire}, dont le r\^ole (tant au point de vue
r\'edactionnel que de celui du travail de mise au point
math\'ematique) est all\'e croissant d'ann\'ee en ann\'ee, le
nom de chaque participant figure en t\^ete des expos\'es dont il
est responsable comme conf\'erencier ou comme r\'edacteur, et la
liste de ceux qui figurent dans un volume d\'etermin\'e se trouve
indiqu\'ee
sur la page de garde dudit volume.

Il convient de donner quelques pr\'ecisions sur les rapports entre
le \emph{Séminaire} et les \emph{Éléments}. Ces derniers
\'etaient destin\'es en principe \`a donner un expos\'e
d'ensemble des notions et techniques jug\'ees les plus fondamentales
dans la G\'eom\'etrie Alg\'ebrique, \`a mesure que ces notions
et techniques elles-m\^emes se d\'egagent, par le jeu naturel
d'exigences de coh\'erence logique et esth\'etique. Dans cette
optique, il \'etait naturel de consid\'erer le
\emph{Séminaire} comme une version pr\'eliminaire des
\emph{Éléments}, destin\'ee \`a \^etre englob\'ee \`a
peu pr\`es totalement, t\^ot ou tard, dans ces derniers. Ce
processus avait d\'ej\`a commenc\'e dans une certaine mesure il
y a quelques ann\'ees, puisque les expos\'es I \`a IV du
pr\'esent volume SGA~1 sont enti\`erement englob\'es par EGA~IV,
et que les expos\'es VI \`a VIII devaient l'\^etre d'ici
quelques ann\'ees dans EGA~VI. Cependant, \`a mesure que se
d\'eveloppe le travail d'\'edification entrepris dans les
\emph{Éléments} et dans le \emph{Séminaire}, et que les
proportions d'ensemble se pr\'ecisent, le principe initial
(d'apr\`es lequel le \emph{Séminaire} ne constituerait qu'une
version pr\'eliminaire et provisoire) appara\^it de moins en
moins r\'ealiste en raison (entre autres) des limites impos\'ees
par la pr\'evoyante nature \`a la dur\'ee de la vie
humaine. Compte tenu du soin g\'en\'eralement apport\'e dans la
r\'edaction des diff\'erentes parties du \emph{Séminaire}, il
n'y aura lieu sans doute de reprendre une telle partie dans les
\emph{Éléments} (ou des trait\'es qui en prendraient la
rel\`eve) que lorsque des progr\`es ult\'erieurs \`a la
r\'edaction permettront d'y apporter des am\'eliorations tr\`es
substantielles, aux prix de modifications assez profondes. C'est le
cas d\`es \`a pr\'esent pour le pr\'esent s\'eminaire SGA~1,
comme on l'a dit plus haut, et \'egalement pour SGA~2 (gr\^ace aux
r\'esultats r\'ecents de Mme~\textsc{Raynaud}). Par contre, rien n'indique
actuellement qu'il en sera ainsi dans un proche avenir pour aucune des
parties cit\'ees plus haut SGA~3 \`a SGA~7.

Les r\'ef\'erences \`a l'int\'erieur du \og S\'eminaire de
G\'eom\'etrie Alg\'ebrique du
Bois Marie\fg sont donn\'ees
ainsi. Une \emph{référence intérieure} \`a une des
parties SGA~1 \`a SGA~7 du S\'eminaire est donn\'ee dans le
style III~9.7, o\`u le chiffre III d\'esigne le num\'ero de
l'expos\'e, qui figure en haut de chaque page de l'expos\'e en
question, et~9.7 le num\'ero de l'\'enonc\'e (ou de la
d\'efinition, remarque, etc.) \`a l'int\'erieur de
l'expos\'e. Le cas \'ech\'eant, des nombres d\'ecimaux plus
longs peuvent \^etre utilis\'es, par exemple 9.7.1, 9.7.2 pour
d\'esigner par exemple les diverses \'etapes dans la
d\'emonstration d'une proposition~9.7. La r\'ef\'erence III~9
d\'esigne la paragraphe 9 de l'expos\'e~III. Le num\'ero de
l'expos\'e est omis pour les r\'ef\'erences internes \`a un
expos\'e. Pour une \emph{référence à une autre} des
\emph{parties} du \emph{Séminaire}, on utilise les m\^emes
sigles, mais pr\'ec\'ed\'es de la mention de la partie en
question des SGA, SGA~1~III~9.7. De m\^eme, la r\'ef\'erence
EGA~IV~11.5.7 signifie: \'El\'ements de G\'eom\'etrie
Alg\'ebrique, Chap.\, IV, \'enonc\'e (ou d\'efinition
etc.)~11.5.7; ici, le premier chiffre arabe d\'esigne encore le
num\'ero du paragraphe. \`A part ces conventions en vigueur dans
l'ensemble des SGA, la bibliographie relative \`a un expos\'e sera
g\'en\'eralement rassembl\'ee \`a la fin de celui-ci, et il y
sera r\'ef\'er\'e \`a l'int\'erieur de l'expos\'e par des
num\'eros entre crochets comme
[\textbf{3}], suivant l'usage.

Enfin, pour la commodit\'e du lecteur, chaque fois que cela semblera
n\'ecessaire, nous joindrons \`a la fin des volumes des SGA un
index des notations, et un index terminologique contenant s'il y a
lieu une traduction anglaise des termes fran\c cais utilis\'es.

Je tiens \`a joindre \`a cette introduction un commentaire
extra-math\'ematique. Au mois de novembre 1969 j'ai eu connaissance
du fait que l'Institut des Hautes \'Etudes Scientifiques, dont j'ai
\'et\'e professeur essentiellement depuis sa fondation, recevait
depuis trois ans des subventions du Minist\`ere des
Arm\'ees. D\'ej\`a comme chercheur d\'ebutant j'ai trouv\'e
extr\^emement regrettable le peu de scrupules que se font la plupart
des scientifiques pour accepter de collaborer sous une forme ou une
autre avec les appareils militaires. Mes motivations \`a ce moment
\'etaient essentiellement de nature morale, donc peu susceptibles
d'\^etres prises au s\'erieux. Aujourd'hui elles acqui\`erent
une force et une dimension nouvelle, vu le danger de destruction de
l'esp\`ece humaine dont nous menace la prolif\'eration des
appareils militaires et des moyens de destruction massives dont ces
appareils disposent. Je me suis expliqu\'e ailleurs de fa\c con
plus d\'etaill\'ee sur ces questions, beaucoup plus importantes
que l'avancement de n'importe quelle science (y compris la
math\'ematique); on pourra par exemple consulter \`a ce sujet
l'article de G.\, Edwards dans le \no 1 du journal Survivre (Ao\^ut
1970), r\'esumant un expos\'e plus d\'etaill\'e de ces
questions que j'avais fait ailleurs. Ainsi, je me suis trouv\'e
travailler pendant trois ans dans une institution alors qu'elle participait \`a
mon insu \`a un mode de financement que je consid\`ere comme
immoral et dangereux\footnote{Il va de soi que l'opinion que je viens
d'exprimer n'engage que ma propre responsabilit\'e, et non pas celle
de la maison d'\'edition Springer qui \'edite le pr\'esent
volume.}. \'Etant \`a pr\'esent seul \`a avoir cette opinion parmi
mes coll\`egues \`a l'IHES, ce qui a condamn\'e \`a
l'\'echec mes efforts pour obtenir la suppression des subventions
militaires du budget de l'IHES, j'ai pris la d\'ecision qui
s'imposait et quitte l'IHES le 30 septembre 1970 et suspends
\'egalement toute collaboration scientifique avec cette institution,
aussi longtemps qu'elle continuera \`a accepter de telles
subventions. J'ai demand\'e \`a M.\, Motchane, directeur de l'IHES,
que l'IHES s'abstienne \`a partir du 1er octobre 1970 de diffuser des textes
math\'ematiques dont je suis auteur, ou faisant partie du
S\'eminaire de G\'eom\'etrie Alg\'ebrique du Bois Marie. Comme
il a \'et\'e dit plus haut, la diffusion de ce s\'eminaire va
\^etre assur\'ee par la maison Julius Springer, dans le s\'erie
des Lecture Notes. Je suis heureux de remercier ici la maison Springer
et Monsieur K.\, Peters pour l'aide efficace et courtoise qu'ils m'ont
apport\'ee pour rendre possible cette publication, en se chargeant
en particulier de la frappe pour la photooffset des nouveaux
expos\'es rajout\'es aux anciens s\'eminaires, et des
expos\'es manquants des s\'eminaires incomplets.

Je remercie \'egalement M.\, J.P.\, Delale, qui s'est charg\'e du
travail ingrat de compiler l'index des notations et l'index
terminologique.

\bigskip\hfill Massy, ao\^ut 1970.

\cleardoublepage
\chapter*{Avertissement}
\thispagestyle{empty}
\label{I.avertissement}


Chacun des expos\'es r\'edig\'es donne la substance de plusieurs
expos\'es oraux cons\'ecutifs. Il n'a pas sembl\'e utile d'en
pr\'eciser les dates.

L'expos\'e~VII, auquel il est r\'ef\'er\'e \`a diverses
reprises au cours de l'expos\'e~\ref{VIII}, n'a pas \'et\'e
r\'edig\'e par le conf\'erencier, qui dans les conf\'erences
orales s'\'etait born\'e \`a esquisser le langage de la descente
dans les cat\'egories g\'en\'erales, en se pla\c cant \`a un
point de vue strictement utilitaire et sans entrer dans les
difficult\'es logiques soulev\'ees par ce langage. Il est apparu
qu'un expos\'e correct de ce langage sortirait des limites des
pr\'esentes notes, ne serait-ce que par sa longueur. Pour un
expos\'e en forme de la th\'eorie de la descente, je renvoie \`a
un article en pr\'eparation de Jean \textsc{Giraud}. En attendant sa
parution\footnote{\label{footnotegiraud}Actuellement paru: J.\, \textsc{Giraud},
\emph{Méthodes de la descente}, M\'emoire \no 2 de la
Soci\'et\'e math\'ematiques de France, 1964.}, je pense qu'un
lecteur attentif n'aura pas de peine \`a suppl\'eer par ses
propres moyens aux r\'ef\'erences fant\^omes de
l'Expos\'e~VIII.

D'autres expos\'es oraux, se pla\c cant apr\`es l'Expos\'e~\ref{XI},
et auxquels il est fait allusion \`a certains endroits du texte,
n'ont pas non plus \'et\'e r\'edig\'es, et \'etaient
destin\'es \`a former la substance d'un Expos\'e~\ref{XII} et d'un
Expos\'e~\ref{XIII}. Les premiers de ces expos\'es oraux reprenaient,
dans le cadre des sch\'emas et des espaces analytiques avec
\'el\'ements nilpotents (tels qu'ils sont introduits dans le
S\'eminaire Cartan 1960/61) la construction de l'espace analytique associ\'e
\`a un pr\'esch\'ema localement de type fini sur un corps
valu\'e complet~$k$, les th\'eor\`emes du type GAGA dans le cas
o\`u $k$ est le corps des complexes, et l'application \`a la
comparaison du groupe fondamental d\'efini par voie transcendante et
le groupe fondamental \'etudi\'e dans ces notes (comparer
A.\, Grothendieck, Fondements de la G\'eom\'etrie Alg\'ebrique,
S\'eminaire Bourbaki \no 190, page~10, d\'ecembre 1959). Les
derniers expos\'es oraux esquissaient la g\'en\'eralisation des
m\'ethodes d\'evelopp\'ees dans le texte pour l'\'etude des
rev\^etements admettant de la ramification mod\'er\'ee, et de la
structure du groupe fondamental d'une courbe compl\`ete priv\'ee
d'un nombre fini de points (comparer \loccit, \no 182, page 27,
th\'eor\`eme~14). Ces expos\'es n'introduisent aucune id\'ee
essentiellement nouvelle, c'est pourquoi il n'a pas sembl\'e
indispensable d'en donner une r\'edaction en forme avant la parution
des chapitres correspondants des \'El\'ements de G\'eom\'etrie
Alg\'ebrique\footnote{Ils sont inclus dans le pr\'esent volume
dans l'Exp.\, \ref{XII} de Mme~Raynaud avec une d\'emonstration
diff\'erente de la d\'emonstration originale expos\'ee dans le
S\'eminaire oral (\cf introduction).}.

Par contre, les th\'eor\`emes du type Lefschetz pour le groupe
fondamental et le groupe de Picard, tant au point de vue local que
global, ont fait l'objet d'un S\'eminaire s\'epar\'e en 1962,
qui a \'et\'e compl\`etement r\'edig\'e et est \`a la
disposition des usagers\footnote{\emph{Cohomologie étale des
faisceaux cohérents et théorèmes de Lefschetz locaux et
globaux} (SGA~2), paru dans North Holland Pub.\ Cie.}. Signalons que
les r\'esultats d\'evelopp\'es tant dans le pr\'esent
S\'eminaire que dans celui de 1962 seront utilis\'es de fa\c con
essentielle dans la parution de plusieurs r\'esultats clefs dans la
cohomologie \'etale des pr\'esch\'emas, qui feront l'objet d'un
S\'eminaire (conduit par M.\, Artin et moi-m\^eme) en 1963/64,
actuellement en pr\'eparation\footnote{\emph{Cohomologie étale
des schémas} (cit\'e SGA~4), \`a para\^itre dans cette
m\^eme s\'erie.}.

Les expos\'es \ref{I} \`a~\ref{IV}, de nature essentiellement locale et
tr\`es \'el\'ementaire, seront absorb\'es enti\`erement par
le chapitre IV des \emph{Éléments de Géométrie
Algébrique}, dont la premi\`ere partie est \`a l'impression et
qui sera sans doute publi\'e vers fin~64. Ils pourront n\'eanmoins
\^etre utiles \`a un lecteur qui d\'esirerait se mettre au
courant des propri\'et\'es essentielles des morphismes lisses,
\'etales ou plats, avant d'entrer dans les arcanes d'un trait\'e
syst\'ematique. Quant aux autres expos\'es, ils seront
absorb\'es dans le chapitre VIII\footnote{En fait, par suite de
modification du plan initialement pr\'evu pour les
\emph{Éléments}, l'\'etude du groupe fondamental y est
report\'ee \`a un chapitre ult\'erieur \`a celui qu'on vient
d'indiquer. Comparer l'introduction qui pr\'ec\`ede le pr\'esent
\og Avertissement\fg.} des \og \'El\'ements\fg, dont la publication ne
pourra gu\`ere \^etre envisag\'ee avant plusieurs ann\'ees.

\bigskip\hfill Bures, juin 1963.





